\documentclass[a4paper,12pt,twoside,BCOR=10mm]{scrbook}

% Packages
\usepackage{ucs}
\usepackage[utf8x]{inputenc}
\usepackage[icelandic, english]{babel}
\usepackage{t1enc}
\usepackage{graphicx}
\usepackage[intoc]{nomencl}
\usepackage{enumerate,color}
\usepackage{url}
\usepackage[pdfborder={0 0 0}]{hyperref}
\usepackage{appendix}
\usepackage{eso-pic}
\usepackage{amsmath}
\usepackage{amsthm}
\usepackage{amssymb}
\usepackage[nottoc]{tocbibind}
\usepackage[sort&compress,authoryear]{natbib}
\usepackage[sf,normalsize]{subfigure}
\usepackage[format=plain,labelformat=simple,labelsep=colon]{caption}
\usepackage{placeins}
\usepackage{tabularx}
\usepackage{mathabx}
\usepackage{mathtools}
% Configurations
\graphicspath{{figs/}}

\newtheorem{theorem}{Theorem}
\newtheorem{corollary}{Corollary}
\newtheorem{remark}{Remark}
\newtheorem{example}{Example}
\newtheorem{lemma}{Lemma}
\newtheorem{definition}{Definition}
\renewcommand{\Re}{\text{Re }}

\setlength{\parskip}{\baselineskip}
\setlength{\parindent}{0cm}
\raggedbottom
% \setkomafont{subsection}{\normalfont\sffamily}

% Eins og templatið á að vera
% \setkomafont{captionlabel}{\itshape}
% \setkomafont{caption}{\itshape}

% Mun fallegri lausn
\setkomafont{captionlabel}{\itshape}
\setkomafont{caption}{\itshape}
\setkomafont{section}{\FloatBarrier\Large}
\setcapwidth[l]{\textwidth}
\setcapindent{1em}


% Times new roman
%\usepackage[T1]{fontenc}
%\usepackage{mathptmx}

%%%%%%%%%%% MODIFY THESE LINES ONLY %%%%%%%%%%%%%%%%%%%%%%%%%%%%%%%%%%%%%%%%%%%%%%%%%%%%%%%%%
\def\thesisyear{2020}       						% Year thesis submitted
\def\thesismonth{XXmonth}						% Month thesis submitted
\def\thesisauthor{Bergur Snorrason}					% Thesis authoreiningaraðferðinni
\def\thesistitle{XXTitle} 						% Title of thesis
\def\thesisshorttitle{XXShort title (50 characters including spaces)} 	% Title of thesis
\def\thesiscredits{XX} 							% Credits awarded for the project
\def\thesissubject{XX}
\def\thesiskind{M.Sc.}							% Masters of PhD thesis
\def\thesiskindformal{Magister Scientiarum}				% Masters of PhD thesis
\def\thesisnroftutors{1}						% Number of tutors
\def\thesisschool{School of Engineering and {Natural Sciences}}		% School
\def\thesisfaculty{XX}							% Faculty
\def\thesisaddress{XXFaculty street address}				% Office address
\def\thesispostalcode{XXFaculty postal code, Reykjavik}			% Office address
\def\thesistelephone{525 4000}						% Office telephone
\def\thesispublisher{XX}						% Publisher
\def\thesistutors{XXNN1 \\ XXNN2}
\def\thesisrepresentative{XXNN3}					% Tutors name
\def\thesiscommittee{XXNN4 \\ XXNN5 }
\def\thesiskeywords{Keyword1, Keyword2, Keyword3}			% Keywords
\def\thesisISBN{XX}           						% Thesis ISBN number
\def\thesisdedication{Dedication}
\def\thesisPrinting{Háskólaprent, Fálkagata 2, 107 Reykjavík}

% Function to add footer to frontpage
\newcommand\BackgroundPic{
\put(0,0){
\parbox[b][\paperheight]{\paperwidth}{
\vfill
\centering
\hspace*{-0.6cm}
\includegraphics[width=\paperwidth,height=\paperheight,
keepaspectratio]{foot}
}}
\setlength{\unitlength}{\paperwidth}
\begin{picture}(0,0)(0,-0.15)
\put(0,0){\color{white}\parbox{1\paperwidth}{\centering\bfseries\sffamily \Large Faculty of \thesisfaculty \\
University of Iceland\\
\thesisyear}}
\end{picture}
}

\begin{document}

\begin{titlepage}
\thispagestyle{empty}
\AddToShipoutPicture*{\BackgroundPic}
%
\begin{center}
\vspace*{1cm}
\includegraphics[width=43.6mm]{ui_1_cmyk}\\
\vspace*{3.0cm}
\huge \sffamily \bfseries \thesistitle

\vspace*{5.5cm}
\normalfont \Large \sffamily \thesisauthor
\AddToShipoutPicture*{\BackgroundPic}
\vfill

\end{center}

\newpage 
\thispagestyle{empty} \mbox{}
\newpage
\vspace*{1.35cm}
\thispagestyle{empty}
\begin{center}

\Large \textbf{\sffamily{\MakeUppercase{\thesistitle}}} \\

\vspace*{1.7cm}
\sffamily{\thesisauthor} \\
\vspace*{1.8cm}
\normalsize \thesiscredits~ECTS thesis submitted in partial fulfillment of a \\
\textit{\thesiskindformal} degree in \thesissubject

\vspace*{1.0cm}
\large
\ifnum\thesisnroftutors >1 Advisors \\ \thesistutors \\ \vspace*{0.4cm}
\else Advisor \\ \thesistutors \\ \vspace*{1.04cm}
\fi
Faculty Representative \\
\thesisrepresentative

\vspace*{0.4cm}
M.Sc. committee \\
\thesiscommittee

\vfill
Faculty of \thesisfaculty \\
\thesisschool \\
University of Iceland \\
Reykjavik, \thesismonth~\thesisyear
\newpage
\end{center}
 \newpage
 \thispagestyle{empty}
 \mbox{} \vfill
 % \setcounter{page}{0} \renewcommand{\baselinestretch}{1.5}\normalsize
 \sffamily{\thesistitle} \\
 \sffamily{\thesisshorttitle} \\
 \thesiscredits ~ECTS thesis submitted in partial fulfillment of a \thesiskind~degree in \thesissubject
\\ \\
Copyright \textcopyright~\thesisyear~ \thesisauthor \\
All rights reserved \\


Faculty of \thesisfaculty \\
\thesisschool \\
University of Iceland \\
\thesisaddress \\
\thesispostalcode, Reykjavik \\
Iceland

Telephone: \thesistelephone \\ \\
\vspace*{\lineskip}

Bibliographic information: \\
\thesisauthor, \thesisyear, \thesistitle, \thesiskind~thesis, Faculty of \thesisfaculty, University of Iceland. \\

ISBN~\thesisISBN

Printing: \thesisPrinting \\
Reykjavik, Iceland, \thesismonth~\thesisyear \\
\newpage
\thispagestyle{empty} \mbox{}
\vfill
\begin{center}\textit{\thesisdedication}\end{center} \vspace*{5cm}
\vfill 

\thispagestyle{empty}
\cleardoublepage
\end{titlepage}


% \dedication{\textit{Dedication} \small \\ Tileinkun má sleppa og skal þá fjarlægja blaðsíðuna. \\
% Tileinkun skal birtast á oddatölu blaðsíðu (hægri síðu).}
\pagenumbering{roman}

\setcounter{page}{5}
\section*{\huge Abstract}
Útdráttur á ensku sem er að hámarki 250 orð.
\vfill \vspace*{1cm}
\section*{\huge Útdráttur}
Hér kemur útdráttur á íslensku sem er að hámarki 250 orð. Reynið að koma útdráttum á eina blaðsíðu en ef tvær blaðsíður eru nauðsynlegar á seinni blaðsíða útdráttar að hefjast á oddatölusíðu (hægri síðu).
\vfill
\newpage

\chapter*{Preface}
Formála má sleppa og skal þá fjarlægja þessa blaðsíðu. Formáli skal hefjast á oddatölu blaðsíðu og nota skal Section Break (Odd Page).

Ekki birtist blaðsíðutal á þessum fyrstu síðum ritgerðarinnar en blaðsíðurnar teljast með og hafa áhrif á blaðsíðutal sem birtist með rómverskum tölum fyrst á efnisyfirliti.

\tableofcontents
\listoffigures
\listoftables



\chapter*{Abbreviations}
\addcontentsline{toc}{chapter}{Abbreviations}
Í þessum kafla mega koma fram listar yfir skammstafanir og/eða breytuheiti. Gefið kaflanum nafn við hæfi, t.d. Skammstafanir eða Breytuheiti. Þessum kafla má sleppa ef hans er ekki þörf. \\
The section could be titled: Glossary, Variable Names, etc.



\chapter*{Acknowledgments}
\addcontentsline{toc}{chapter}{Acknowledgments}
Í þessum kafla koma fram þakkir til þeirra sem hafa styrkt rannsóknina með fjárframlögum, aðstöðu eða vinnu. T.d. styrktarsjóðir, fyrirtæki, leiðbeinendur, og aðrir aðilar sem hafa á einhvern hátt aðstoðað við gerð verkefnisins, þ.m.t. vinir og fjölskylda ef við á. Þakkir byrja á oddatölusíðu (hægri síðu).

\pagenumbering{arabic}
\setcounter{page}{1}
\chapter{Introduction}



\chapter{Preliminaries}
\section{Measure theory}
\section{Functional analysis}
\subsection{Hahn-Banach}
\subsection{Riesz representation theorem}



\chapter{Rudin-Carleson theorem}
\begin{theorem}
Let $E$ be a closed subset of $\mathbb{T}$ of Lebesgue-measure $0$, let $f$ be a continuous function on $E$ and let $T$ be a simply connected subset of $\mathbb{C}$ such that $f(\mathbb{\overline{D}}) \subset T$.
Then there exists an $F \in \mathcal{A}$, such that $F = f$ on $E$ and $F(\mathbb{\overline{D}}) \subset T$. (TODO skilgreina allt)
\end{theorem}
We will break the proof into several (TODO how many?) lemmas.
\begin{lemma}
Let $E$ be a clsoed set of Lebesbue-measure $0$.
Then there exists an integrable function $\mu > 1$ such that
	$\mu$ is continuous on $\mathbb{T} \setminus E$,
	$\mu = +\infty$ on $E$,
	if $w \in E$ then $\mu(z) \xrightarrow{z \rightarrow w} +\infty$, and
	$\mu$ has a bounded derivative on any closed subarc of $\mathbb{T} \setminus E$.
\end{lemma}
\begin{proof}
The function $\mu$ is found by solving the Dirichlet problem. (TODO Thomas Ransford p. 95 4.2.6) (TODO finish this)
\end{proof}
\begin{lemma}
If $f$ is a simple continuous function on $E$ such that $\Re f \geq 0$, then there exists an $F \in \mathcal{A}$ such that $F = f$ on $E$ and $\Re F \geq 0$ on $\overline{\mathbb{D}}$.
\end{lemma}
\begin{proof}
We will show that this holds in the case where $E = E_0 \cup E_1$, $f = 0$ on $E_0$, $f = \alpha \neq 0$ on $E_1$ and $\Re \alpha \geq 0$.
This suffices since simple functions are finite linear combinations of characteristic functions.
\end{proof}
\begin{proof}
TODO
\end{proof}
\begin{corollary}[Fatou]
Let $E$ be a closed subset of $\mathbb{T}$ of Lebesgue-measure $0$.
There exists a function $f \in \mathcal{A}$ that vanishes on $E$ and nowhere else.
\end{corollary}
\begin{proof}
It's clear from the theorem (TODO add ref) that there exists a function $f \in \mathcal{A}$ that vanishes on $E$.
TODO
\end{proof}




\chapter{F. and M. Riesz theorem}
In this section we will endeavour to show that the annihilating measures of $\mathcal{A}|_{\mathbb{T}}$ are absolutely continuous with respect to the Lebesgue measure.
We will show this to be a corollary of the F. and M. Riesz theorem, which we will prove in the manner of Rudin. %TODO add the reference to Rudin real and complex analysis.
To attain the main result of this section we need some lemmas and definitions. %TODO rewrite these two lines
To prove one of the lemmas we will also use the following famous theorem:
\begin{definition}
% TODO finish these
Pointwise bounded and equicontinuity.
\end{definition}
\begin{theorem}[Ascoli-Arzela] %TODO prove this or move it to preliminaries
Let $\mathcal{F}$ be a pointwise bounded equicontinuous collection of complex functions on a metric space $(X, d)$, and $X$ contains a countable dense subset.  %TODO change to 'X is separable'
Then every sequence in $\mathcal{F}$ contains a subsequence that converges uniformly on every compact subsets of $X$.
\end{theorem}
\begin{proof}
Let $E$ be a countable dense subset of $X$,
	$\{f_n\}$ be a series in $\mathcal{F}$,
	and $x_1, x_2, ...$ be an enumeration of $E$.
\end{proof}
\begin{definition}
% TODO finish these
Poisson kernel, Poisson integral, Poisson integral of a measure.
\end{definition}
\begin{lemma}
\label{FMRlemma1}
Let $\mu$ be a complex Borel measure, and $u = P[d\mu]$.
Then
\[
	\|u_r\|_1 = \frac{1}{2\pi} \int_{-\pi}^{\pi} |u(re^{i\theta})| d\theta \leq \|\mu\|.
\]
\end{lemma}
\begin{proof}
TODO
\end{proof}
\begin{lemma}
\label{FMRlemma2}
Let $f \in H^1$.
Then there exists a $g \in L^1(\mathbb{T})$ such that $f = P[g]$.
\end{lemma}
\begin{proof}
TODO
\end{proof}
\begin{lemma}
%11.30(a) in Rudin
\label{FMRlemma3}
Let $u$ be harmonic in $\mathbb{D}$ and
\[
	\sup_{0 < r < 1} \|u_r\|_1 = M < \infty. %TODO define u_r
\]
Then there exists a unique complex Borel measure $\mu$ on $\mathbb{T}$ such that $u = P[d\mu]$.
\end{lemma}
We will need the following lemma in the proof of \ref{FMRlemma3}:
\begin{lemma}
\label{FMRlemma31}
Let $X$ be a separable Banach space,
	$\{ \Gamma_n\}$ be a sequence of linear functionals on $X$,
	and $\sup_n \|\Gamma_n\| = M < \infty$.
Then there exists a subsequence $\{\Gamma_{n_i}\}$ such that the limit 
\[
	\Gamma x = \lim_{i \rightarrow \infty} \Gamma_{n_i} x
\]
exists for every $x \in X$.
We also have that $\Gamma$ is linear and $\|\Gamma\| \leq M$.
\end{lemma}
\begin{proof}
\end{proof}
\begin{proof}[Proof of \ref{FMRlemma3}]
% TODO fix the Riesz rep thm ref. It's 6.19 in Rudin.
\end{proof}

\begin{theorem}[F. and M. Riesz theorem]
If $\mu$ is a complex Borel measure on $\mathbb{T}$ and
\[
	\int e^{-int} d\mu = 0
\]
for $n = -1, -2, ...$, then $\mu \lll m$.
\end{theorem}
\begin{proof}
Let $f = P[d\mu]$.
If we set $z = re^{i\theta}$ we get that
\[
	P(z, e^{it})
		= P_r(\theta - t)
		= \sum_{n \in \mathbb{Z}} r^{|n|}e^{in(\theta - t)}
		= \sum_{n \in \mathbb{Z}} r^{|n|}e^{in\theta}e^{-int}.
\]
We can use the assumption of the theorem to write $f$ as a power series by
\begin{align*}
% TODO Fubini?
	f(z) 
		&= \int_{\mathbb{T}} P(z, e^{it}) d\mu(e^{it}) \\
		&= \int_{\mathbb{T}} \sum_{n \in \mathbb{Z}} r^{|n|}e^{in\theta}e^{-int} d\mu(e^{it}) \\
		&= \sum_{n \in \mathbb{Z}} r^{|n|}e^{in\theta} \int_{\mathbb{T}} e^{-int} d\mu(e^{it}) \\
		&= \sum_{n = 0}^{\infty} r^n e^{in\theta} \int_{\mathbb{T}} e^{-int} d\mu(e^{it}) \\
		&= \sum_{n = 0}^{\infty} \hat{\mu}_n z^n,
\end{align*}
where $\hat{\mu}_n$ is the $n$-th Fourier coefficient of $\mu$.
This along with \ref{FMRlemma1} gives us that $f \in H^1$. % TODO make this more clear
We can now define a $g \in H^1$, by \ref{FMRlemma2}, such that $f = P[g]$.
It follows from \ref{FMRlemma3} that $d\mu = f d\sigma$. %TODO define \sigma as Lebesgue on T scaled to 1.
TODO
\end{proof}
\begin{corollary}
Let $A$ be the closed subspace of $C(\mathbb{T})$ that consists of all functions that are restriction from $\mathcal{A}$ (TODO makes this def global and maybe call it something else).
All measures in $A^{\bot}$ are absolutely continuous with regards to the Lebesgue-measure on $\mathbb{T}$ (TODO define all).
\end{corollary}
\begin{proof}
Let $\mu \in A^{\bot}$. By definition we have that
\[
	\int fd\mu = 0.
\]
Now since $t \mapsto e^{-int}$ is entire for $n = -1, -2, ...$ we have that their restriction to $\mathbb{T}$ are in $A$.
Thus,
\[
	\int e^{-int} d\mu = 0
\]
for all $n = -1, -2, ...$ and $\mu \lll m$.
\end{proof}



\chapter{A generalization of the Rudin-Carleson theorem}
This borrows from Bishop (reference TODO).
\begin{theorem}[General Rudin-Carleson theorem]
Let $X$ be a compact Hausdorff space,
	$V = (C(X), \| \cdot \|_{\infty})$,
	$B$ be a closed subspace of $C(X)$,
	$B^{\bot}$ be the annihilating measures for $B$,
	$S$ be a closed subset of $X$,
	and $f$ be a continues function on $S$.
If $\int_S f d\mu = 0$ holds for all $\mu \in B^{\bot}$ then there exists a function $F \in B$ such that $F = f$ on $S$.
\end{theorem}
\begin{proof}
Since $f$ is continuous and $S$ is a closed subset of a compact set, and therefore also compact, $f$ is bounded.
So we can, with out loss of generality, assume that $|f| < r < 1$ on $S$.
Let $U_r$ be the subset of $B$ defined by $U_r = \{g;\ \|g\| < r\}$ and $\phi$ be the mapping from $B$ to $C(S)$ that sends a member of $B$ to its restriction on $S$.
It suffices to show that $f \in \phi(U_r)$.
Let's first show that $f \in \overline{\phi(U_r)} =: V_r$, by assuming otherwise, and showing it leads to a contradiction.

We now assume $f \not \in V_r$.
By Hahn-Banach (TODO ref) we can define a bounded linear functional $\alpha$, such that $\alpha(f) > 1$ and $|\alpha(h)| < 1$, for $h \in V_r$.
We can then define a measure $\mu_1$ by the Riesz-representation theorem (TODO ref) that fulfills
\[
	\alpha(g) = \int g d\mu_1
\]
for all $g \in C(S)$.
We will refer to the associated functional on $B$ by $\beta(g) = \phi(\alpha(g))$.
Since $\phi(g) \in V_r$ for all $g \in U_r$ we have that
\[
	\beta(g) = \alpha(\phi(g)) < 1,
\]
for all $g \in U_r$,
	due to the construction of $\alpha$.
From this we get
\begin{align*}
	\| \beta \|
	&= \sup \{ |\beta(g)|;\ |g| < 1 \}\\
	&= \sup \{ (1/r)|\beta(g)|;\ |g| < r \}\\
	&\leq \sup \{ (1/r);\ |g| < r \}\\
	&= 1/r.
\end{align*}
Let's denote the Riesz representation of $\beta$ by $\mu_2$, set $\mu = \mu_1 - \mu_2$ and see that $\mu \in B^{\bot}$.
But
\[
	0 = \left | \int_S f d\mu \right | 
		\geq \int_S f d\mu_1 - r\|\mu_2\| 
		\geq \int_S f d\mu_1 - r\frac{1}{r}
		> 1 - r \frac{1}{r} = 0,
\]
where the first equality is the assumption in the theorem.
This is the contradiction that gives that $f \in V_r$.
We can now take a $F_1$ in $U_r$, and therefore also in $B$ such that $|f - F_1| < \lambda/2$ on $S$, with $\lambda := 1 - r$.
Remember that $F_1 \ in U_r$ implies that $\|F_1\| < r$.
Now let $f_1 = f - F_1$ and use the same method as above to obtain an $F_2$ such that $\|F_2\| < \lambda/2$ and $|f - F_2| < \lambda/4$ on $S$.
Iterating this process yields a series $(F_n)_{n \in \mathbb{N}}$ from $B$ that fulfill $\|F_n\| < 2^{1 - n}\lambda$ for $n > 1$ and
\[
	\left | f - \sum_{k = 1}^n F_k \right | < 2^{-n}\lambda
\]
on $S$ for $n > 1$.
We finally let 
\[
F = \sum_{k = 1}^{\infty} F_k.
\]
Now $F \in B$,
\[
	\|F\| \leq \|F_1\| + \|F - F_1\| = r + \sum_{k = 2}^{\infty}2^{1 - n}\lambda = r + \lambda = 1,
\]
and $F = f$ on $S$. (TODO bæta við matinu, svo línan að ofan meiki sens)
\end{proof}
\begin{corollary}
Let $X$ be a compact Hausdorff space,
	$V = (C(X), \| \cdot \|_{\infty})$,
	$B$ be a closed subspace of $C(X)$,
	$B^{\bot}$ be the annihilating measures for $B$,
	$S$ be a closed subset of $X$,
	and $f$ be a continues function on $S$.
If $S$ is $B^{\bot}$-null (TODO define this) then there exists a function $F \in B$ such that $F = f$ on $S$.
\end{corollary}
\begin{proof}
If $S$ is $B^{\bot}$-null we have that $\int_S f d\mu = 0$ for all $\mu \in B^{\bot}$ (TODO add ref to above).
\end{proof}
\begin{remark}
The corollary is the version of the theorem from Bishop (TODO add ref here).
Note also that if we set $X = \mathbb{T}$ and $B = \mathcal{A}$ we can use F. and M. Riesz (TODO add ref here) to prove the classical Rudin-Carleson theorem (TODO add ref here).
\end{remark}
It is of course worth noting an applications of (TODO REF to theorem) where the corollary fails.
\begin{example}
Let $X = \mathbb{T}$,
	$B = \mathcal{A}$ (TODO define this),
	$E$ be a closed $m$-null (TODO define this, an maybe change to $m_{\sigma}$) subset of $\partial \mathbb{T}$ that is not dense in $E$,
	$F = \{ e^{i\theta};\ a \leq \theta \leq b \}$,
	and choose $a$ and such that $E$ and $F$ are disjoint and $a \neq b$.
The last assumption restricts us to $E$ that are not dense in the $\mathbb{E}$.
Since $a \neq b$ we obtain that $S := E \cup F$ does not fulfill the requirements of the Rudin-Carelson theorem (TODO add ref here) nor the above corollary (TODO add ref).
Let's choose $f$ such that $f = 0$ on $F$, and $f$ is continues on $S$.
We now have for all $\mu \in \mathcal{A}^{\bot}$
\begin{align*}
	\left | \int_S f d\mu \right |
	&= \left |\int_E f d\mu + \int_F f d\mu \right |\\
	&\leq \left |\int_E f d\mu \right | + \left | \int_F f d\mu \right |\\
	&= 0 + \left | \int_F f d\mu \right |\\
	&= 0.
\end{align*}
The F. and M. Riesz theorem (TODO add ref here) tells us the since $E$ is $m$-null it is also $\mu$-null, which gives the third step.
The final step stems from the fact that $f$ vanishes on $F$.
We now see the $X$, $B$, and $f$ are all as in theorem (TODO add ref) so there exists a $F \in B$, such that $F = f$ on S.
\end{example}

\appendix
\renewcommand{\chaptername}{Appendix}
\chapter{Annað}

\end{document}
